% === START SEGMENT 2 ===

%% ----------------------------------------------------------------
%%  SECTION III — SYSTEM ARCHITECTURE
%% ----------------------------------------------------------------
\section{System Architecture}
\label{sec:architecture}

This section presents the end-to-end architecture of RobustIDPS, which
integrates a multi-branch surrogate intrusion detection ensemble with a
continual-learning--reinforcement-learning (CL-RL) unified framework and an
autonomous response pipeline.  Fig.~\ref{fig:architecture} provides a
high-level overview of the system.

\begin{figure}[t]
\centering
\fbox{\parbox{0.95\columnwidth}{\vspace{4cm}\centering\textit{System Architecture Diagram}}}
\caption{Overall architecture of RobustIDPS showing the 7-branch surrogate
ensemble, CL-RL unified framework, and autonomous response pipeline.}
\label{fig:architecture}
\end{figure}

\subsection{SurrogateIDS: 7-Branch Ensemble}
\label{sec:surrogate}

The detection backbone, termed \textit{SurrogateIDS}, comprises seven
parallel branch networks whose predictions are fused to yield both a
classification decision and a calibrated uncertainty estimate.  Each branch
$b_j$, $j\in\{1,\dots,7\}$, receives the same 83-dimensional input feature
vector $\mathbf{x}\in\mathbb{R}^{83}$ extracted from raw network traffic and
processes it through three fully connected hidden layers of dimensions
$[512,\,256,\,128]$.  Every hidden layer is followed by Batch
Normalisation~\cite{ioffe2015batch}, a ReLU activation, and Dropout with
probability $p=0.3$:
\begin{equation}
  \mathbf{h}_j^{(\ell)} = \mathrm{Dropout}\!\bigl(
    \mathrm{ReLU}\!\bigl(
      \mathrm{BN}\!\bigl(
        \mathbf{W}_j^{(\ell)}\,\mathbf{h}_j^{(\ell-1)}
        + \mathbf{b}_j^{(\ell)}
      \bigr)
    \bigr),\; p
  \bigr),
  \label{eq:branch}
\end{equation}
where $\mathbf{h}_j^{(0)}=\mathbf{x}$.  The seven branch outputs
$\{\mathbf{h}_j^{(3)}\}_{j=1}^{7}$, each of dimension~128, are
concatenated to form a 896-dimensional representation that is projected
through a shared fusion layer:
\begin{equation}
  \mathbf{z} = \mathrm{ReLU}\!\bigl(
    \mathbf{W}_{\mathrm{fuse}}\,
    [\mathbf{h}_1^{(3)};\,\dots;\,\mathbf{h}_7^{(3)}]
    + \mathbf{b}_{\mathrm{fuse}}
  \bigr),\quad
  \mathbf{W}_{\mathrm{fuse}}\in\mathbb{R}^{256\times 896},
  \label{eq:fusion}
\end{equation}
yielding $\mathbf{z}\in\mathbb{R}^{256}$.  A final linear classifier maps
$\mathbf{z}$ to 34 output classes (one benign class and 33 distinct attack
categories):
\begin{equation}
  \hat{\mathbf{y}} = \mathrm{softmax}\!\bigl(
    \mathbf{W}_c\,\mathbf{z} + \mathbf{b}_c
  \bigr),\quad
  \mathbf{W}_c\in\mathbb{R}^{34\times 256}.
  \label{eq:classifier}
\end{equation}

\subsection{MC-Dropout Uncertainty Quantification}
\label{sec:mcdropout}

At inference time we retain dropout in stochastic mode and perform
$T=20$ forward passes through the ensemble.  Let
$\hat{\mathbf{y}}^{(t)}$ denote the softmax output of the $t$-th pass.
The mean predictive distribution is
$\bar{\mathbf{y}}=\frac{1}{T}\sum_{t=1}^{T}\hat{\mathbf{y}}^{(t)}$.
We quantify epistemic uncertainty via \emph{predictive entropy}
\begin{equation}
  \mathcal{H}[\bar{\mathbf{y}}]
  = -\sum_{c=1}^{34}\bar{y}_c\,\log\bar{y}_c,
  \label{eq:pred_entropy}
\end{equation}
and \emph{mutual information}
\begin{equation}
  \mathcal{I}
  = \mathcal{H}[\bar{\mathbf{y}}]
    - \frac{1}{T}\sum_{t=1}^{T}\mathcal{H}[\hat{\mathbf{y}}^{(t)}].
  \label{eq:mi}
\end{equation}
Samples whose mutual information exceeds a calibrated threshold are
flagged as out-of-distribution, triggering the continual learning module
described in Section~\ref{sec:ewc}.

\subsection{SDE-TGNN for Temporal Network Dynamics}
\label{sec:sde_tgnn}

To capture the continuous-time evolution of host-level communication
graphs we employ a Stochastic Differential Equation--Temporal Graph
Neural Network (SDE-TGNN).  The latent state
$\mathbf{h}(t)\in\mathbb{R}^{d}$ of each node evolves according to the
It\^{o} SDE
\begin{equation}
  \mathrm{d}\mathbf{h}
  = f_\theta(\mathbf{h},t)\,\mathrm{d}t
    + g_\phi(\mathbf{h},t)\,\mathrm{d}\mathbf{W},
  \label{eq:sde}
\end{equation}
where $f_\theta\colon\mathbb{R}^{d}\times\mathbb{R}\to\mathbb{R}^{d}$
is a drift network, $g_\phi\colon\mathbb{R}^{d}\times\mathbb{R}\to
\mathbb{R}^{d\times d}$ is a diffusion network, and $\mathbf{W}$ is a
standard Wiener process.  The graph attention mechanism aggregates
neighbour states at each integration step, enabling the model to detect
subtle lateral-movement patterns that unfold across irregular time
intervals.


%% ----------------------------------------------------------------
%%  SECTION IV — METHODOLOGY
%% ----------------------------------------------------------------
\section{Methodology}
\label{sec:methodology}

\subsection{Continual Learning with Elastic Weight Consolidation}
\label{sec:ewc}

Network intrusion landscapes evolve over time as adversaries introduce
novel attack variants.  To mitigate catastrophic forgetting while
accommodating new traffic distributions, we adopt Elastic Weight
Consolidation (EWC)~\cite{kirkpatrick2017overcoming} augmented with
experience replay.

\subsubsection{EWC Loss Formulation}
When training on a new task~$\mathcal{T}_t$, the overall loss is
\begin{equation}
  \mathcal{L}(\theta)
  = \mathcal{L}_{\mathrm{CE}}(\theta)
    + \frac{\lambda}{2}\sum_{k}F_k\bigl(\theta_k - \theta_k^{*}\bigr)^{2},
  \label{eq:ewc}
\end{equation}
where $\mathcal{L}_{\mathrm{CE}}$ is the cross-entropy loss on the
current task, $\theta^{*}$ are the optimal parameters from the preceding
task, $F_k$ is the $k$-th diagonal element of the Fisher Information
Matrix (FIM), and $\lambda$ controls the consolidation strength.

\subsubsection{Experience Replay}
We maintain a fixed-capacity replay buffer
$\mathcal{M}$ of size $|\mathcal{M}|=5000$ samples, updated via
reservoir sampling~\cite{vitter1985random} so that each previously
observed class retains approximately equal representation.  At every
training step a mini-batch is drawn uniformly from $\mathcal{M}$ and
mixed with the current-task batch to compute
$\mathcal{L}_{\mathrm{CE}}$.

\subsubsection{Online Fisher Estimation}
Computing the full FIM after each task is prohibitively expensive.
Instead, we maintain an exponential moving average:
\begin{equation}
  \hat{F}^{(t)}
  = \alpha\,\hat{F}^{(t-1)} + (1-\alpha)\,\tilde{F}^{(t)},
  \quad \alpha=0.5,
  \label{eq:fisher_ema}
\end{equation}
where $\tilde{F}^{(t)}$ is the empirical Fisher computed on a
representative subset of task-$t$ data.  This formulation smoothly
discounts older consolidation signals and is well suited to the
streaming setting of network monitoring.

\begin{algorithm}[t]
\DontPrintSemicolon
\SetAlgoLined
\SetKwInOut{Input}{Input}\SetKwInOut{Output}{Output}
\Input{Current task data $\mathcal{D}_t$; previous parameters
$\theta^{*}$; Fisher estimate $\hat{F}^{(t-1)}$; replay buffer
$\mathcal{M}$; EWC weight $\lambda$; EMA coefficient $\alpha$}
\Output{Updated parameters $\theta$; updated $\hat{F}^{(t)}$;
updated $\mathcal{M}$}
Initialise $\theta \leftarrow \theta^{*}$\;
\For{each mini-batch $(\mathbf{X}_t,\mathbf{y}_t)$ from
$\mathcal{D}_t$}{
  Sample replay batch $(\mathbf{X}_r,\mathbf{y}_r)\sim\mathcal{M}$\;
  $\mathcal{L}_{\mathrm{CE}}\leftarrow
    \mathrm{CE}(\theta;\mathbf{X}_t,\mathbf{y}_t)
    + \mathrm{CE}(\theta;\mathbf{X}_r,\mathbf{y}_r)$\;
  $\mathcal{L}_{\mathrm{EWC}}\leftarrow
    \frac{\lambda}{2}\sum_k \hat{F}_k^{(t-1)}
    (\theta_k-\theta_k^{*})^2$\;
  $\mathcal{L}\leftarrow
    \mathcal{L}_{\mathrm{CE}}+\mathcal{L}_{\mathrm{EWC}}$\;
  $\theta\leftarrow\theta-\eta\,\nabla_\theta\mathcal{L}$\;
  Update $\mathcal{M}$ via reservoir sampling with
  $(\mathbf{X}_t,\mathbf{y}_t)$\;
}
Compute empirical Fisher $\tilde{F}^{(t)}$ on subset of
$\mathcal{D}_t$\;
$\hat{F}^{(t)}\leftarrow\alpha\,\hat{F}^{(t-1)}
  +(1-\alpha)\,\tilde{F}^{(t)}$\;
$\theta^{*}\leftarrow\theta$\;
\Return $\theta,\;\hat{F}^{(t)},\;\mathcal{M}$\;
\caption{EWC with Experience Replay and Online Fisher Estimation}
\label{alg:ewc}
\end{algorithm}

\subsection{Constrained RL Response Agent}
\label{sec:cpo}

Once the SurrogateIDS ensemble produces a prediction and uncertainty
estimate, an autonomous response agent selects a defensive action.  We
formalise this as a Constrained Markov Decision Process (CMDP).

\subsubsection{CMDP Formulation}
The CMDP is defined by the tuple
$(\mathcal{S},\mathcal{A},P,R,C,d_0,\gamma)$, where $\mathcal{S}$ is
the state space, $\mathcal{A}$ is the action space, $P$ denotes
transition dynamics, $R$ is the reward function, $C$ is a cost function,
$d_0$ is the initial state distribution, and $\gamma=0.99$ is the
discount factor.

\paragraph{State space.}
The state $\mathbf{s}\in\mathbb{R}^{55}$ is a concatenation of: (i)~the
34-dimensional softmax output $\hat{\mathbf{y}}$, (ii)~the 7-branch
pre-fusion embeddings reduced to 14~dimensions via PCA, (iii)~predictive
entropy, mutual information, and five contextual features (traffic rate,
session duration, protocol type encoding, source reputation score,
historical alert count), and (iv)~the previous action one-hot encoding
of dimension~5, yielding $34+14+7+5=55$ dimensions in total, after accounting for uncertainty signals within the contextual block.

\paragraph{Action space.}
The discrete action set is
$\mathcal{A}=\{\textsc{Monitor},\,\textsc{RateLimit},\,
\textsc{Reset},\,\textsc{Block},\,\textsc{Quarantine}\}$
with associated severity costs
$\mathbf{c}=[0,\,0.5,\,1,\,2,\,5]$.

\subsubsection{Constrained Policy Optimisation (CPO)}
The agent maximises expected cumulative reward subject to a safety
constraint on cumulative cost:
\begin{align}
  \max_{\pi}\;&\;
    \mathbb{E}_{\pi}\!\Bigl[\sum_{t=0}^{\infty}\gamma^t\,
    R(s_t,a_t)\Bigr], \label{eq:cpo_obj}\\
  \text{s.t.}\;&\;
    \mathbb{E}_{\pi}\!\Bigl[\sum_{t=0}^{\infty}\gamma^t\,
    C(s_t,a_t)\Bigr] \le \varepsilon. \label{eq:cpo_constraint}
\end{align}
The constraint threshold $\varepsilon$ is calibrated so that the
false-positive blocking rate remains below $0.1\%$ on the validation
set.  We solve~\eqref{eq:cpo_obj}--\eqref{eq:cpo_constraint} using the
CPO algorithm~\cite{achiam2017constrained}, which performs trust-region
updates on a surrogate objective while projecting the policy back into
the feasible set at each iteration.

\paragraph{Policy network.}
The policy $\pi_\psi$ is parameterised by a two-hidden-layer network
$[256,\,128]$ with ReLU activations, mapping the 55-dimensional state to
a categorical distribution over five actions:
\begin{equation}
  \pi_\psi(a\mid s)
  = \mathrm{softmax}\!\bigl(
    \mathbf{W}_2\,\mathrm{ReLU}(
      \mathbf{W}_1\,\mathbf{s}+\mathbf{b}_1)
    +\mathbf{b}_2
  \bigr).
  \label{eq:policy}
\end{equation}

\begin{algorithm}[t]
\DontPrintSemicolon
\SetAlgoLined
\SetKwInOut{Input}{Input}\SetKwInOut{Output}{Output}
\Input{Initial policy $\pi_\psi$; value network $V_\omega$; cost
value network $V_\omega^C$; constraint threshold $\varepsilon$;
trust-region radius $\delta$; discount $\gamma$}
\Output{Trained policy $\pi_\psi$}
\For{episode $= 1,2,\dots$}{
  Collect trajectory $\tau=\{(s_t,a_t,r_t,c_t)\}$ under $\pi_\psi$\;
  Compute advantages $\hat{A}_t$ and cost advantages $\hat{A}_t^C$
  via GAE($\lambda$)\;
  Estimate constraint value
  $J_C(\pi_\psi)\approx\frac{1}{|\tau|}
    \sum_t\gamma^t c_t$\;
  Compute surrogate objective
  $L(\psi)=\frac{1}{|\tau|}\sum_t
    \frac{\pi_\psi(a_t|s_t)}{\pi_{\psi_{\mathrm{old}}}(a_t|s_t)}
    \hat{A}_t$\;
  Compute surrogate cost
  $L_C(\psi)=\frac{1}{|\tau|}\sum_t
    \frac{\pi_\psi(a_t|s_t)}{\pi_{\psi_{\mathrm{old}}}(a_t|s_t)}
    \hat{A}_t^C$\;
  \eIf{$J_C(\pi_\psi)\le\varepsilon$}{
    $\psi\leftarrow\arg\max_{\psi'}L(\psi')$ s.t.\
    $\overline{D}_{\mathrm{KL}}(\pi_{\psi_{\mathrm{old}}}\|
      \pi_{\psi'})\le\delta$\;
  }{
    Project: $\psi\leftarrow\arg\max_{\psi'}L(\psi')$ s.t.\
    $L_C(\psi')\le\varepsilon - J_C$ and
    $\overline{D}_{\mathrm{KL}}\le\delta$\;
  }
  Update $V_\omega,\;V_\omega^C$ by regression on returns\;
}
\Return $\pi_\psi$\;
\caption{Constrained Policy Optimisation (CPO) for Response Agent}
\label{alg:cpo}
\end{algorithm}

\subsection{Unified Fisher Information Framework}
\label{sec:unified_fisher}

A central contribution of RobustIDPS is the \emph{unified Fisher
information framework} that couples the detection and response modules
through a single regularisation signal.  Let $\hat{F}_k^{\mathrm{det}}$
denote the diagonal Fisher element for parameter~$k$ of the
SurrogateIDS ensemble (Section~\ref{sec:ewc}), and let
$[F_\pi]_{kk}$ denote the corresponding diagonal element of
the policy Fisher for parameters shared between the two modules.  We
define the unified Fisher estimate as
\begin{equation}
  \hat{F}_k^{\mathrm{unified}}
  = \beta\,\hat{F}_k^{\mathrm{det}}
    + (1-\beta)\,[F_\pi]_{kk},
  \quad \beta=0.7.
  \label{eq:unified_fisher}
\end{equation}
The weighting $\beta=0.7$ biases consolidation toward preserving
detection accuracy---the primary objective of an IDPS---while still
allowing sufficient plasticity for the RL policy to adapt its response
strategy to evolving threat landscapes.  Equation~\eqref{eq:unified_fisher}
replaces $F_k$ in~\eqref{eq:ewc} during joint training, thereby
preventing catastrophic interference between the detection and response
objectives.

\subsection{Adversarial Robustness}
\label{sec:adversarial}

We evaluate and harden RobustIDPS against six representative
adversarial attack strategies that span gradient-based, optimisation-based,
geometry-based, and stochastic perturbation categories:
\begin{enumerate}
  \item \textbf{FGSM}~\cite{goodfellow2015explaining}: single-step
    perturbation with $\varepsilon=0.1$.
  \item \textbf{PGD}~\cite{madry2018towards}: iterative projected
    gradient descent with $\varepsilon=0.1$, step size $\alpha=0.01$,
    and 40 iterations.
  \item \textbf{C\&W}~\cite{carlini2017towards}: $\ell_2$ optimisation
    attack with confidence parameter $\kappa=0$ and 50 iterations.
  \item \textbf{DeepFool}~\cite{moosavi2016deepfool}: minimal-norm
    perturbation with a maximum of 30 iterations.
  \item \textbf{Gaussian noise}: additive isotropic noise with standard
    deviation $\sigma=0.1$.
  \item \textbf{Label masking}: stochastic label corruption with flip
    probability $p_{\mathrm{flip}}=0.1$, simulating poisoning attacks on
    the training pipeline.
\end{enumerate}

\subsubsection{Adversarial Training Procedure}
During each training epoch we generate adversarial examples on-the-fly
using PGD (the strongest first-order attack in our suite) and augment
the training batch:
\begin{equation}
  \tilde{\mathbf{x}} = \Pi_{\mathcal{B}_\varepsilon(\mathbf{x})}
  \!\bigl(\mathbf{x} + \alpha\,\mathrm{sign}(
    \nabla_{\mathbf{x}}\mathcal{L}_{\mathrm{CE}}(\theta;\mathbf{x},y)
  )\bigr),
  \label{eq:pgd_step}
\end{equation}
where $\Pi_{\mathcal{B}_\varepsilon}$ denotes projection onto the
$\ell_\infty$ ball of radius $\varepsilon$.  The loss is computed over
both clean and adversarial samples with equal weighting, following the
min-max formulation of~\cite{madry2018towards}:
\begin{equation}
  \min_\theta\;\mathbb{E}_{(\mathbf{x},y)\sim\mathcal{D}}
  \Bigl[\max_{\|\boldsymbol{\delta}\|_\infty\le\varepsilon}
    \mathcal{L}_{\mathrm{CE}}(\theta;\mathbf{x}+\boldsymbol{\delta},y)
  \Bigr].
  \label{eq:adv_train}
\end{equation}

\subsection{Byzantine-Resilient Federated Learning}
\label{sec:federated}

In realistic multi-site deployments, raw traffic data cannot be
centralised due to privacy regulations.  RobustIDPS therefore supports
federated training across $N$ participating sites while tolerating up to
$f<N/3$ Byzantine nodes.

\subsubsection{Aggregation Strategies}
We implement and compare three aggregation protocols:
\begin{itemize}
  \item \textbf{FedAvg}~\cite{mcmahan2017communication}: vanilla
    weighted averaging of client updates.
  \item \textbf{FedProx}~\cite{li2020federated}: FedAvg augmented with
    a proximal regularisation term weighted by $\mu=0.01$ to handle
    statistical heterogeneity.
  \item \textbf{FedGTD}: a game-theoretic aggregation scheme in which
    client contributions are weighted by a trust score derived from
    historical validation performance.
\end{itemize}

\subsubsection{Byzantine Detection}
Before aggregation, the server projects all client gradient updates into
a low-dimensional subspace of dimension $d_{\mathrm{proj}}=10$ and
computes pairwise cosine similarities.  Any client whose mean similarity
to the remaining clients falls below the threshold $\tau=0.3$ is flagged
as potentially Byzantine and excluded from the current round.  The
surviving gradients are combined via coordinate-wise trimmed mean,
discarding the highest and lowest values at each coordinate.

\subsubsection{Differential Privacy}
To provide formal privacy guarantees, each client clips its gradient
update to norm $S$ and adds calibrated Gaussian noise:
\begin{equation}
  \tilde{\mathbf{g}}_i
  = \mathrm{clip}(\mathbf{g}_i,\,S)
    + \mathcal{N}\!\bigl(\mathbf{0},\,\sigma_{\mathrm{DP}}^2 S^2
    \mathbf{I}\bigr),
  \label{eq:dp}
\end{equation}
achieving $(\epsilon_{\mathrm{DP}},\delta)$-differential privacy per
round via the moments accountant~\cite{abadi2016deep}.

\subsubsection{Stochastic Game Dynamics and Nash Equilibrium}
We model the interaction between the federated defender and an adaptive
adversary as a continuous-time stochastic game.  Let
$\mathbf{v}(t)\in\mathbb{R}^{m}$ denote the joint strategy profile.
Its dynamics are governed by an SDE with Poisson jumps:
\begin{equation}
  \mathrm{d}\mathbf{v}
  = \mu(\mathbf{v},t)\,\mathrm{d}t
    + \sigma(\mathbf{v},t)\,\mathrm{d}\mathbf{W}
    + \int_{\mathbb{R}^m}\!\mathbf{J}(\mathbf{v},\mathbf{z})\,
      \widetilde{N}(\mathrm{d}t,\mathrm{d}\mathbf{z}),
  \label{eq:jump_sde}
\end{equation}
where $\widetilde{N}$ is a compensated Poisson random measure capturing
sudden strategy shifts (e.g., new zero-day exploits).  Under mild
Lipschitz conditions on $\mu$ and $\sigma$, the strategy profile
converges to a Nash equilibrium in the mean-square sense, ensuring that
the federated defence remains robust even against strategically adaptive
adversaries.  A formal convergence proof is provided in the
supplementary material.

% === END SEGMENT 2 ===
